\section{Unitary representations and their star algebras}
\label{sec:characters_unitary}

\begin{definition}[Star algebra of a unitary representation]
    \label{def:peps_unitary_representation_algebra}
    \lean{Representation.euclideanMatrixRepresentation,
        Representation.generatedMatrixStarAlgebra,
        Representation.matrixInvariantSubrepresentation}
    \leanok
    Let $U:G\to\mathrm U(D)$ be a unitary representation. Regard its
    matrices as operators on the standard Hilbert space $\C^D$,
    and let $\mathcal A_U$ be their generated unital star algebra.
    Every $\mathcal A_U$-invariant subspace carries the restriction of $U$
    \cite[Section~4.1]{Schuch2010PEPS}.
\end{definition}

\begin{lemma}[Group and algebra invariant subspaces]
    \label{thm:peps_unitary_representation_algebra}
    \lean{Subrepresentation.isAtom_of_isIrreducible_toRepresentation,
        Representation.isInvariantSubspace_generatedMatrixStarAlgebra_iff,
        Representation.isIrreducibleSubspace_generatedMatrixStarAlgebra_iff,
        Representation.isIntertwiner_generatedMatrixStarAlgebra_iff,
        Representation.sameIsotype_generatedMatrixStarAlgebra_iff}
    \leanok
    \uses{def:peps_unitary_representation_algebra}
    A subspace is invariant under $\mathcal A_U$ if and only if it is
    invariant under every $U(g)$. Its restricted group representation
    is irreducible if and only if it is an irreducible invariant subspace
    of $\mathcal A_U$. An operator intertwines the algebra on an invariant
    subspace if and only if it intertwines every group operator there.
    Two irreducible invariant subspaces have the same algebra isotype
    if and only if their restricted group representations are equivalent.
\end{lemma}

\begin{proof}\leanok
    Since $U(g)^*=U(g^{-1})$, closure under adjoints introduces no new
    generators. Invariance and intertwining extend from the generators
    by sums, products, and scalar multiples. Thus the two notions have
    identical invariant subspaces. Irreducibility is minimality among
    nonzero invariant subspaces, or equivalently the atom condition for
    the lattice of subrepresentations. A nonzero intertwiner between
    irreducible subspaces is an isomorphism; conversely, extending an
    isomorphism of the restricted representations to the ambient space
    by first projecting orthogonally onto its source subspace produces an
    algebra intertwiner.
\end{proof}

\begin{lemma}[Orthonormal irreducible matrix blocks]
    \label{thm:peps_unitary_irreducible_blocks}
    \lean{Representation.exists_unitary_irreducible_matrix_blocks}
    \leanok
    \uses{thm:peps_unitary_representation_algebra}
    Suppose $G$ is finite. There are positive integers $d_i,m_i$ and an
    orthonormal basis $(e_{i,a,j})$, with $1\leq a\leq m_i$ and
    $1\leq j\leq d_i$, such that every row spans an irreducible
    subrepresentation. Rows belonging to distinct indices $i$ have
    distinct characters. For each $g\in G$ there are matrices $D_i(g)$
    satisfying
    \begin{align}
        U(g)e_{i,a,j}=\sum_{j'}D_i(g)_{j'j}e_{i,a,j'}
        \label{eq:peps_unitary_irreducible_rows}
    \end{align}
    \cite[Sections~4.1 and~7]{Schuch2010PEPS}.
\end{lemma}

\begin{proof}\leanok
    Applying the orthonormal adapted-basis decomposition of a
    finite-dimensional star algebra to $\mathcal A_U$, the preceding
    correspondence identifies its irreducible rows and distinct isotypes with
    irreducible group representations and distinct characters, and its common
    matrix action on the multiplicity rows gives
    (\ref{eq:peps_unitary_irreducible_rows}).
\end{proof}

\begin{lemma}[Character multiplicities of orthonormal rows]
    \label{thm:peps_unitary_character_blocks}
    \lean{Representation.exists_unitary_character_matrix_blocks}
    \leanok
    \uses{thm:peps_unitary_irreducible_blocks, thm:peps_isotypic_projectors}
    In the decomposition (\ref{eq:peps_unitary_irreducible_rows}), each
    row has dimension $d_i$, and the character of a row occurs in the
    full representation with multiplicity exactly $m_i$.
    Thus the orthonormal decomposition gives the dimensions and
    multiplicities appearing in \cite[Lemma~4.4]{Schuch2010PEPS}.
\end{lemma}

\begin{proof}\leanok
    Map the basis of one row to another row of the same component.
    The common matrix action makes this a group intertwiner, so those
    rows have the same character. Rows of different components have
    distinct characters. The character projector onto component $i$
    therefore fixes exactly its $m_id_i$ orthonormal basis vectors.
    Its trace is $m_id_i$. The character trace identity gives the same
    trace as $m_{\chi_i}d_i$. Since $d_i>0$, cancellation gives
    $m_{\chi_i}=m_i$.
\end{proof}

\begin{lemma}[Regular representations in relative coordinates]%
    \label{thm:peps_regular_blocking_coordinates}
    \lean{Matrix.permMatrixHom_prodCongr,
        RegularRepresentation.blockingEquiv, RegularRepresentation.blockingEquiv_apply,
        RegularRepresentation.blockingEquiv_symm_apply,
        RegularRepresentation.blockingEquiv_conjugates,
        RegularRepresentation.blockingFamilyEquiv,
        RegularRepresentation.blockingFamilyEquiv_apply,
        RegularRepresentation.blockingFamilyEquiv_symm_apply,
        RegularRepresentation.blockingFamilyEquiv_smul,
        RegularRepresentation.blockingFamilyEquiv_conjugates,
        RegularRepresentation.blockingMatrix,
        RegularRepresentation.blockingMatrix_mem_unitaryGroup,
        RegularRepresentation.blockingMatrix_conjugates,
        RegularRepresentation.blockingFamilyMatrix,
        RegularRepresentation.blockingFamilyMatrix_mem_unitaryGroup,
        RegularRepresentation.blockingFamilyMatrix_conjugates}
    \leanok
    The basis permutation $T\ket{a,b}=\ket{a,ab}$ is unitary and satisfies
    $T^\dagger(L_g\otimes L_g)T=L_g\otimes\Id$, where
    $L_g\ket a=\ket{ga}$. For a finite family of additional labels, the
    permutation $\ket{a,(b_i)_i}\mapsto\ket{a,(ab_i)_i}$ similarly transforms
    simultaneous left multiplication into left multiplication of $a$
    alone \cite[Section~6]{Schuch2010PEPS}.
\end{lemma}

\begin{proof}\leanok
    The inverse change of coordinates is $(a,(c_i)_i)\mapsto
        (a,(a^{-1}c_i)_i)$. Thus simultaneous multiplication by $g$ sends
    $(a,(ab_i)_i)$ to $(ga,(gab_i)_i)$, and in relative coordinates only
    the first label changes. The permutation matrices are unitary, and
    the matrix of a product permutation is the Kronecker product of the
    factor matrices.
\end{proof}
